% Part: incompleteness
% Chapter: theories-computability
% Section: interpretability

\documentclass[../../../include/open-logic-section]{subfiles}

\begin{document}

\olfileid{inc}{tcp}{itp}

\olsection{النظريات التي تقبل تأويل $\Th{Q}$ غير قابلة للقرار}

يمكننا تعزيز هذه النتائج أكثر. على نحو غير صوري، ينطوي تأويل لغة
$\Lang{L_1}$ في لغة أخرى~$\Lang{L_2}$ على تعريف الكون ورموز العلاقات
ورموز الدوال في~$\Lang{L_1}$ بواسطة !!{formula}s في~$\Lang{L_2}$.
ومع أننا لن نستغرق الوقت في ذلك، يمكن ضبط هذا التعريف بدقة.

\begin{thm}
  افترض أن $\Th{T}$ نظرية في لغة يمكن فيها تأويل لغة الحساب على نحو
  تكون معه $\Th{T}$ متسقة مع تأويل~$\Th{Q}$. عندئذ تكون $\Th{T}$ غير
  قابلة للقرار. وإذا أثبتت $\Th{T}$ تأويل بديهيات~$\Th{Q}$، فلا يكون
  أي امتداد متسق لـ$\Th{T}$ قابلًا للقرار.
\end{thm}

ليس البرهان إلا تعديلًا يسيرًا لبرهان المبرهنة السابقة؛ إذ يمكن
استعمال مثال مضاد للحصول على فصل بين $\Th{Q}$ و$\Th{\bar Q}$. ويمكن
اتخاذ $\Th{ZFC}$، أي نظرية مجموعات زرميلو--فرينكل مع بديهية الاختيار،
أساسًا بديهيًا قويًا بما يكفي لإجراء قدر كبير من الرياضيات المألوفة.
ويمكن في~$\Th{ZFC}$ تعريف الأعداد الطبيعية، وتصدق بديهيات~$\Th{Q}$
وفق هذا التأويل. ولذلك لدينا:

\begin{cor}
لا يوجد امتداد متسق قابل للقرار لـ$\Th{ZFC}$.
\end{cor}

\begin{cor}
لا توجد نظرية تامة متسقة !!{axiomatizable} حسابيًا تمتدّ~$\Th{ZFC}$.
\end{cor}

ليس في لغة~$\Th{ZFC}$ إلا علاقة ثنائية واحدة هي~$\in$. (بل لا تحتاج
في الواقع حتى إلى الهوية.) ولذلك لدينا:

\begin{cor}
منطق الرتبة الأولى لأي لغة فيها رمز علاقة ثنائية غير قابل للقرار.
\end{cor}

تمتد هذه النتيجة إلى كل لغة فيها رمزا دالة أحاديان، إذ يمكن استعمالهما
لمحاكاة رمز علاقة ثنائية. والنتائج المذكورة للتو مضبوطة: فقد تبين أن
منطق الرتبة الأولى للغة ليس فيها إلا رموز علاقات \emph{أحادية} ورمز
دالة \emph{أحادي} واحد على الأكثر قابل للقرار.

وثمة معلومة طريفة أخرى. نعلم أن مجموعة الجمل الصادقة في النموذج
القياسي من اللغة ذات الرموز $\Obj 0$، و$'$، و$+$، و$\times$، و$<$ غير
قابلة للقرار. ويمكن في الواقع تعريف~$<$ بدلالة الرموز الأخرى، ثم تعريف
$+$ بدلالة $\times$ و$'$. ومن ثم تكون مجموعة الجمل الصادقة في اللغة
ذات الرموز $\Obj 0$، و$'$، و$\times$ غير قابلة للقرار. ومن جهة أخرى،
أثبت بريسبورغر أن مجموعة الجمل في اللغة ذات الرموز $\Obj 0$، و$'$،
و$+$ الصادقة في النموذج القياسي قابلة للقرار. غير أن الإجراء غير قابل
للتنفيذ حسابيًا من الناحية العملية.

\end{document}
